\chapter{Keeping the encoder and decoder in agreement}
\MARSChapterHeading{Keeping the encoder and decoder in agreement}
\label{revision:h263-reconstruction}

A video encoder can predict a region of the current picture from a picture it has
already sent. If the prediction is close, the difference may be cheaper to encode
than the original samples. This creates a dependency: the next prediction uses
what the decoder reconstructed from the previous compressed picture. The encoder
must retain a corresponding reconstruction of its own.

This excerpt assumes that you know the capstone's picture planes, motion
compensation, eight-by-eight transform, quantization and coefficient coding.
It develops the connection between those stages. The full capstone produces an
actual H.263 stream using a documented restricted baseline feature set. Its input
is a short, script-generated QCIF sequence: a $176\times144$ luma plane and two
$88\times72$ chroma planes per picture. The input limit is at most 30 pictures
per second and ten seconds. Those bounds make the assignment finite; they do
not change the reconstruction rule.

\section{A reference picture is a computed result}

Suppose an original sample is 100, but its coded representation reconstructs as
98. The next picture's desired sample is 104. An encoder that subtracts its old
original value obtains a residual of four. A decoder that adds four to its stored
98 obtains 102. There is no arithmetic error in either subtraction or addition.
The error is using different reference values.

For this simplified example, using 98 on both sides gives residual six and
reconstruction 104. In the real codec the residual also passes through transform,
quantization and inverse operations. Those operations can change it. The encoder
therefore cannot assume that sending a correction reproduces the original sample
exactly. It must compute the same reconstruction path it expects at the decoder.

Keep two buffers with distinct meanings. The raw input holds the samples we want
to encode. The reconstruction holds samples available for later prediction.
Treating raw input as the reconstruction can make later pictures disagree even
when the first picture looks plausible.

\begin{figure}[htbp]
  \centering
  \begin{tikzpicture}[node distance=7mm,>=Stealth,
      every node/.style={draw,align=center,font=\small,text width=96mm},
      alt={Current samples minus prediction produce residuals. The forward transform and quantizer produce selected levels that are coded into the stream. A local path inverse-quantizes those same levels, applies the inverse transform, adds the same prediction and clips to sample range. The resulting reconstructed picture supplies future predictions.}]
    \node (residual) {Current samples minus prediction\\signed residuals};
    \node[below=of residual] (levels) {Forward transform and quantization\\selected levels};
    \node[below=of levels] (inverse) {Inverse quantization and inverse transform\\reconstructed residuals};
    \node[below=of inverse] (pixels) {Add the same prediction and clip\\reconstructed samples};
    \draw[->] (residual)--(levels);
    \draw[->] (levels)--(inverse);
    \draw[->] (inverse)--(pixels);
  \end{tikzpicture}
  \caption{The encoder's local reconstruction consumes the selected levels. Those levels also supply coefficient coding for the output stream.}
  \label{revision:local-reconstruction}
\end{figure}

\section{Follow the values across representation boundaries}

For an INTER block, let $S[y,x]$ be an unsigned current sample and $P[y,x]$ its
motion-compensated prediction. Form the signed residual
\[
  E[y,x]=S[y,x]-P[y,x].
\]
Both samples lie in $[0,255]$, so $E$ lies in $[-255,255]$. Subtracting in byte
lanes would lose that signed range. Widen before subtraction and retain signed
halfword residuals.

The forward transform expresses the spatial residual block as coefficients.
Quantization selects integer levels from those coefficients. The levels are what
the coefficient coder describes to the decoder; the original coefficients are no
longer available there. Local reconstruction must inverse-quantize the selected
levels and inverse-transform the resulting coefficients. Call its spatial result
$D[y,x]$.

The reconstructed sample is
\[
  X[y,x]=\min(255,\max(0,P[y,x]+D[y,x])).
\]
The addition occurs while $D$ is still signed. For the capstone's inverse boundary,
$D$ lies in $[-256,255]$, so the sum lies in $[-256,510]$ and fits a signed halfword.
Narrow to a byte after clipping. With prediction 250 and residual 20, the exact
sum is 270 and the output is 255; byte wrapping would give 14. With prediction
five and residual $-10$, the exact sum is $-5$ and the output is zero. Clipping
the residual to zero before adding would incorrectly give five.

For an INTRA block there is no inter-picture prediction to add. The inverse
transform describes absolute sample values, which are clipped to $[0,255]$.
The capstone's baseline INTRA path transforms sample values directly. A level
shift copied from a different codec would change this path.

\section{Work through a constant block}

A constant block isolates the reconstruction loop because only its DC transform
coefficient is nonzero in the ideal transform. Take all 64 current samples as
130 and all prediction samples as 128. Every residual is two. With the transform
normalization used here, a constant residual $r$ has DC coefficient $8r$, so the
forward DC is sixteen.

Use quantizer parameter $Q=8$. The capstone's ordinary INTER quantizer selects
level one for this coefficient. Its inverse quantization reconstructs coefficient
23. That is the coefficient to feed to the inverse transform. Using the original
sixteen would undo the quantizer's choice locally even though the decoder must
work from level one.

The ideal inverse of a DC-only block with coefficient 23 is $23/8=2.875$ at every
sample. The supplied tested transform rounds this case to three. Adding that
residual to prediction 128 gives reconstruction 131. The raw input was 130.
The one-unit difference is a consequence of this quantized representation; it
is not by itself evidence that the block pipeline is wrong.

\begin{center}
  \tagpdfsetup{table/header-rows={1}}
  \begin{tabular}{lll}
    \toprule
    Boundary               & Value    & Meaning                               \\
    \midrule
    Current and prediction & 130, 128 & Unsigned samples                      \\
    Transform input        & 2        & Signed residual                       \\
    Forward DC             & 16       & Original coefficient                  \\
    Selected level         & 1        & Value represented in the stream       \\
    Reconstructed DC       & 23       & Inverse-quantized coefficient         \\
    Inverse output         & 3        & Rounded spatial correction            \\
    Reconstruction         & 131      & Sample retained for future prediction \\
    \bottomrule
  \end{tabular}
\end{center}

This is a trace for the named transform and quantizer, not a requirement that
every accurate implementation choose identical values at every rounding boundary.
Use the transform's accuracy contract when comparing implementations. Within one
implementation, the selected levels and the local reconstruction must remain
consistent with the bitstream it actually emits.

\Needspace{14\baselineskip}
\begin{MARSWorkedExample}{The next picture}
  Let the next desired block contain 132 throughout, with zero motion. Prediction
  from the reconstructed 131 gives residual one and forward DC eight. At $Q=8$,
  this falls in the quantizer's zero region: the selected levels are all zero.
  The inverse residual is zero and reconstruction stays at 131. The encoder and
  decoder can agree while both differ from the current raw input.

  Now repeat the calculation using raw 130 as the encoder's reference. It gives
  residual two, selected level one and inverse residual three. The faulty local
  encoder adds three to 130 and obtains 133. The decoder adds three to its actual
  reference 131 and obtains 134. The disagreement is additional to the original
  quantization error. Changing the next raw picture cannot repair the faulty
  ownership of the reference buffer.
\end{MARSWorkedExample}

\section{Keep the vector path on the same boundaries}

MSA can accelerate the sample operations around the transform. A prediction and
an inverse residual can occupy halfword lanes. Add them at that width, apply
signed lower and upper bounds, then pack the clipped values into bytes. The
range proof above justifies ordinary halfword addition for this stage. It does
not justify byte addition, and it does not determine the widths required inside
the transform's dot products.

Treat each vector access as a claim about the image layout. An eight-sample row
occupies eight bytes, while an MSA vector occupies sixteen. A sixteen-byte load
from the start of an arbitrary strided row can read padding or another region
outside the admitted footprint. Gather the active bytes, or prove that a larger
access consists entirely of permitted samples. A large backing allocation does
not supply that proof.

Use an independent scalar calculation to check the vector stage. Compare the
signed inverse residuals as well as the final bytes. Final clipping can hide a
wrong intermediate: both 270 and 300 become 255. A test that inspects only the
clipped picture would miss that difference. Include negative corrections and
values just inside and outside the sample bounds.

\section{Publish the picture that was actually encoded}

A picture-level encoder has more than one output. It produces compressed bytes
and a reconstructed picture that may become the reference for a later picture.
These outputs describe the same coding decisions. Publishing only one leaves
subsequent work without a coherent accepted state.

Stage the picture's coding output and reconstruction while the operation is in
progress. Retain the previous accepted reference until the picture succeeds.
If a child operation fails, the caller must be able to identify which state is
still valid. The capstone's acceptance boundary includes completing the required
output write before advancing its accepted-reference state. A partially written
file cannot in general be rolled back; the driver reports failure and stops
instead of continuing from a picture the receiver does not have completely.

A candidate motion search reads the current picture and the retained reference.
It does not modify either one. A chosen predictor and its reconstruction belong
to the picture being built. Keeping these roles separate allows a failed block,
capacity check or write to leave the last accepted reference usable for diagnosis.

Identical consecutive raw pictures do not imply that a macroblock should be
skipped. The current raw samples may still differ from the preceding reconstructed
reference because of earlier quantization. Evaluate a skip decision against the
prediction and syntax rules that the decoder will use.

\section{Find where agreement first breaks}

Inspect the stages in dependency order: prediction, selected levels,
inverse-quantized coefficients, inverse output and reconstructed samples. A matching
forward transform does not establish a matching local decoder. If the first
reconstructed picture agrees but the next does not, verify the reference identity
before changing the transform. If disagreement appears at the lower sample bound,
look for an unsigned interpretation of a negative correction or clipping performed
too early.

An external H.263 decoder supplies a separate check on the emitted stream.
Use it to verify dimensions, picture count, decodability and the documented quality
conditions. Do not use a host encoder's private motion choices as the required
student answer: several conforming choices can describe the same sequence.
The capstone also checks inverse-transform accuracy and its own internal contracts;
one successful external decode is not sufficient evidence for every component.

\begin{MARSExercise}{Separate loss from disagreement}
  Use the two constant-block traces to write down three quantities for each picture:
  the desired raw sample, the encoder's local reconstruction and the decoder's
  reconstruction. Identify which differences can result from quantization and which
  show that the two sides used different state. Then replace the constant block with
  an asymmetric pattern to expose transform orientation errors hidden by DC-only data.
\end{MARSExercise}

\section*{Sources and capstone connection}
The reconstruction and coding boundaries follow the restricted baseline profile
specified for the capstone and \href{https://www.itu.int/rec/T-REC-H.263-200501-I/en}{ITU-T Recommendation H.263},
including its inverse-transform and reconstructed-sample boundaries. Quantizer
choices and the exact constant-block trace above belong to the documented teaching
implementation. The public \path{docs/h263-profile.md} records the full project
scope and points to the script-generated input and staged lab work. This excerpt
explains the local reconstruction loop; it does not replace the complete encoder,
syntax tables or conformance checks.
